\section{Related Work}\label{sec:related}

\subsection{Re-identification with frozen and foundation encoders}
CLIP-ReID \cite{li2023clipreid} established that a language-image backbone is a strong
initialisation for person ReID, and a family of adaptation methods followed --- prompt
learning, parameter-efficient tuning and instruction conditioning
\cite{he2024mambapro,chen2024instructreid,zhu2022pass}. These works fine-tune. The
complementary question, what a foundation encoder is worth \emph{without} adaptation, is
answered less often and less consistently. A 2026 cross-paradigm study
\cite{paradigm2026} compares supervised, self-supervised and language-aligned models over 11
backbones and 9 datasets and is the closest prior work to ours in scope; a parallel line
evaluates multimodal LLMs as the matcher itself \cite{mmreid2025}. Neither varies the readout
as an experimental axis: each fixes a single evaluation convention and reports backbone
differences under it. Our contribution is orthogonal and, we argue, prior to theirs --- if the
ranking depends on the convention, then a table that fixes one convention is reporting a
property of the pair, not of the backbones.

\subsection{Agglomerative vision foundation models}
AM-RADIO \cite{ranzinger2024amradio} showed that several heterogeneous teachers --- a
language-aligned model, a self-supervised model and a segmentation model --- can be distilled
into a single student that matches or exceeds each teacher on that teacher's own benchmarks,
and RADIOv2.5 \cite{heinrich2025radiov25} addressed the mode-switching and resolution
pathologies that the first version exhibited. The line has since been extended by other
agglomerative students \cite{eupe2026} and by mixture-of-experts variants
\cite{amoe2025}. What this family has not been evaluated on is instance-level retrieval, where
the relevant quantity is not whether a class is separable but whether one individual is
separable from a visually near-identical other. Distillation from an averaged teacher signal
is a plausible threat to exactly that margin, and it is the hypothesis \cref{sec:ablation}
tests against C-RADIOv4's own teachers.

\subsection{Probing a frozen representation}
Linear probing is the standard instrument for comparing frozen representations, inherited from
the self-supervised classification literature \cite{chen2020simclr,caron2021dino,oquab2024dinov2}.
Two conventions travel under the name and are not interchangeable. The classification
convention trains a linear classifier on frozen features and reports \emph{the classifier's}
accuracy; the retrieval convention discards the classifier and reports mAP over the features
it reshaped. Foundation-model reports overwhelmingly use the first, ReID uses the second, and
numbers from one are quoted beside numbers from the other. Within ReID, the readout that
matters in practice is a margin-based one --- ArcFace \cite{deng2019arcface},
CosFace \cite{wang2018cosface} or Circle loss \cite{sun2020circle} --- because it optimises the
angular geometry the retrieval metric reads, which is precisely why substituting it for a plain
linear head is not a neutral implementation choice. We report the retrieval convention
throughout and treat the head as an axis of the experiment rather than as part of its setup.

\subsection{Evaluation protocol and its failure modes}
Market-1501 \cite{zheng2015scalable} fixed the query--gallery--junk convention that most person
benchmarks still follow; MSMT17 \cite{wei2018msmt} scaled it; CUHK03-NP
\cite{zhong2017reranking} supplied the detector-box regime and the 767/700 split that made
CUHK03 comparable across papers. mINP \cite{ye2022deep} was introduced precisely because mAP is
dominated by easy true matches and says little about the hardest one, which is the quantity a
deployed system fails on. Beyond ReID, OpenOOD \cite{zhang2023openood} is the clearest recent
demonstration that protocol discipline --- fixed splits, tuning on validation data only,
stratified difficulty --- changes published conclusions in a whole subfield, and the
open-set ReID literature \cite{liao2014operid,li2018adversarial,wang2020gom} has repeatedly
observed that closed-set rank metrics assume the answer is present in the gallery. Our
substrate (\cref{sec:substrate}) is built on the same premise: an evaluation protocol should
be a value with an identity, so that two numbers can be checked for comparability rather than
assumed to be comparable.

\subsection{Benchmark toolboxes}
Torchreid \cite{zhou2019torchreid} and FastReID \cite{he2023fastreid} are the field's reference
implementations and are training frameworks first; their evaluation code is a component of a
training loop rather than an independently usable artifact, and neither exposes an
agglomerative backbone, a frozen-probe harness, or an open-set protocol. We release the
evaluation half as a standalone package with no training loop, no dataset mirror and no torch
dependency in its core, which is what allows a run record to be re-scored on a machine that
cannot run the encoder that produced it.
